\section{When a sharper collective response destabilises feedback}
\label{hrec:sec:problem}
Consider a detector array that repeatedly revises a common drive. Each detector returns a binary answer to a small signed signal. A positive majority asks for a positive correction, a negative majority for a negative one. Increasing the number of independent detectors makes the collective answer more decisive near the decision threshold. This is useful when the task is to distinguish two fixed alternatives. Once that answer changes the drive that generates the next observation, however, decisiveness becomes a feedback gain. The question is no longer only whether the detectors recognise a bias. It is whether the apparatus settles after they act on it.

A hierarchy changes this construction before it changes any differential equation. We can collect all detector outputs in one vote, or first form local majorities and then vote on their verdicts. Both arrangements use the same leaves. They do not have the same response curve. Nor does adding a level mean that we can move the resulting structure cheaply. An instruction applied at its base may need several updates to reach its top, and a prescribed final configuration can require cancellation of unwanted motion at intermediate levels. Static discrimination, closed-loop stability and the effort of an intervention are three questions about the same architecture. They require three explicit constructions.

This use of hierarchy is architectural rather than an assumption about institutional authority. Simon's definition encompasses interrelated subsystems nested within subsystems, and explicitly distinguishes that broader structure from a formal chain of command \cite{Simon1962236d85}. Feedback through an outcome that an intervention itself changes also has a precise contemporary formulation in performative prediction. There, the deployed predictor changes the distribution on which it is subsequently evaluated, and stability of retraining is distinct from risk optimality \cite{Perdomo2020a5b2ff}. Neither observation determines the gain of our detector array. To obtain that gain we must say what is aggregated, what the controller observes and when it acts.

We do this for two deterministic sampled models. One adjusts a broadcast by the gradient of a target response using an old report. The other first advances a relaxing plant, then applies a collective response to the newly observed mismatch. The placement of the response changes both the stability calculation and the nonlinear result. We then replace a flat detector bank by exact, specified majority trees, and finally attach an independently defined actuator. The contribution is the resulting chain of calculations: from a microscopic vote to its derivative, from that derivative to the appropriate full-loop boundary, and from an admitted architecture to a reachable intervention within budget. The derivations and numerical method are given below.

